\section{Related Work}
\label{sec:related}

\subsection{JEPA, SIGReg, and \LeWM{}}
\label{sec:rel-jepa}

Joint Embedding Predictive Architectures learn predictive
representations by forecasting future embeddings rather than pixels
\cite{assran2023jepa}.
The JEPA objective is typically a distance between a predicted latent
and a target latent from a context encoder, often with stop-gradient
on the target path so the predictor cannot collapse by copying a
degenerate constant.
Video JEPA variants extend this idea to temporal sequences
\cite{bardes2024vjepa}, still avoiding reconstruction decoders that
can dominate loss and encourage shortcut visual features.
A recurring failure mode is \emph{representation collapse}: embeddings
that ignore inputs yet trivially satisfy a prediction objective.
SIGReg~\cite{balestriero2023sigreg} addresses collapse by pushing
projected embedding coordinates toward an isotropic Gaussian via a
sketched Epps--Pulley statistic, without requiring a stop-gradient
teacher network as in contrastive or BYOL-style recipes.

\LeWM{}~\cite{maes2026lewm} brings JEPA to end-to-end pixel world
models for control.
Its recipe is deliberately simple: a ViT encoder trained from scratch,
an action-conditioned predictor, MSE next-embedding loss, and SIGReg
with weight $\lambda{=}0.09$ and $M{=}1024$ projections.
Planning uses CEM in latent space against a goal embedding derived
from a future observation.
Reported PushT success under full-data training is high
($96{\pm}2.83\%$), which motivates careful protocol matching when
studying predictor variants on smaller subsets: a weak absolute CEM
score on 1k windows is expected and should not be confused with a
failed reproduction of the full recipe.
We adopt \LeWM{}'s official encoder, SIGReg recipe, projectors, and CEM
evaluator unchanged, so our study isolates predictor dynamics rather
than re-tuning the full stack.

\subsection{Continuous-time and ODE world models}
\label{sec:rel-ode}

Neural ODEs~\cite{chen2018neuralode} parameterize continuous
transformations via $\mathrm{d}z/\mathrm{d}t = v_\theta(z,t)$ and
backpropagate through ODE solvers.
This viewpoint has influenced continuous-depth networks, continuous
normalizing flows, and---more recently---world models that treat
latent evolution as a flow rather than a discrete jump.

\ODEWorld{}~\cite{liu2026odeworld} introduces Physical-Time Flow: a
continuous latent velocity field integrated by an ODE solver for video
and control, emphasizing dynamical representation structure beyond
discrete next-step prediction.
Our setting is narrower and deliberately comparable: we keep
\LeWM{}'s JEPA loss and planning loop, and redefine a single discrete
planning step as integrate $\tau{:}0\to 1$ under an
action-conditioned velocity field.
We do not claim the full Physical-Time Flow agenda (adaptive physical
time, multi-scale horizons, or ODEWorld's broader architecture);
instead we ask whether the ODE inductive bias alone helps under a
matched \LeWM{} protocol.

\subsection{Discrete AR vs.\ continuous latent dynamics}
\label{sec:rel-disc-cont}

Discrete autoregressive predictors map $(z_t,a_t)\mapsto\hat z_{t+1}$
in one shot.
They are efficient at inference (one forward pass per step) and match
the discrete action schedule used by CEM rollouts.
Continuous predictors instead learn $v_\theta$ and obtain
$\hat z_{t+1}$ by numerical integration.
This costs more compute per step (RK4 with $N{=}4$ evaluates the
velocity field four times per substep interval, sixteen times over
$\tau{:}0\to1$ with $N{=}4$), but exposes intermediate states along
the interval and encourages residual updates ($z(1)\approx z(0)$ when
$v\approx 0$).
Our experiments quantify both sides: loss and CEM on a fixed split,
plus wall-clock CEM cost on the same hardware.

\subsection{Latent MPC and CEM planning}
\label{sec:rel-cem}

Model-based RL and planning systems often learn latent dynamics and
optimize action sequences by shooting
\cite{chua2018pets,hafner2019planet}.
Cross-entropy method (CEM) samples candidate action sequences, rolls
out the learned dynamics, ranks by a goal-conditioned cost, and
refits a sampling distribution around elite samples.
\LeWM{} uses this interface with a latent goal embedding derived from
a future observation at a fixed goal offset.

We use the official \LeWM{} CEM evaluator on a fixed validation window
set so discrete and ODE predictors share the same starts and
feasibility filter (\Cref{sec:setup}).
Earlier probes on random full-dataset starts are discarded as
protocol-mismatched; all planning claims in this paper use the
val-99 / 93-feasible protocol tied to \texttt{validate/loss}.

\subsection{Positioning summary}
\label{sec:rel-position}

Relative to prior work, \method{} is intentionally conservative:
\begin{itemize}[leftmargin=*,itemsep=2pt]
\item vs.\ \LeWM{}: same JEPA + SIGReg + CEM stack; only the predictor
map changes from discrete AR to RK4-integrated $v_\theta$.
\item vs.\ \ODEWorld{}: we borrow the continuous-velocity idea, but do
not adopt Physical-Time Flow, adaptive physical time, or ODEWorld's
full training recipe.
\item vs.\ classical Neural ODEs: we integrate on a fixed normalized
interval $\tau\in[0,1]$ per planning step, with piecewise-constant
actions, rather than learning an unrestricted continuous-time policy.
\end{itemize}
This positioning keeps comparisons fair on official PushT windows and
makes negative or tied CEM results interpretable rather than confounded
by simultaneous encoder or planner changes.
